% !TEX program = pdflatex
% ============================================================================
% Boilerplate LaTeX - Dossier professionnel / Cahier des charges
% Remplace les placeholders entre crochets : [Nom du projet], [Client], etc.
% Compilation conseillée : pdflatex dossier_professionnel_boilerplate.tex
% ============================================================================

\documentclass[11pt,a4paper]{report}

% ----------------------------------------------------------------------------
% Encodage, langue, mise en page
% ----------------------------------------------------------------------------
\usepackage[utf8]{inputenc}
\usepackage[T1]{fontenc}
\usepackage[french]{babel}
\usepackage{lmodern}
\usepackage{microtype}
\usepackage{geometry}
\geometry{margin=2.4cm}
\usepackage{setspace}
\onehalfspacing
\usepackage{parskip}

% ----------------------------------------------------------------------------
% Images, tableaux, liens, couleurs
% ----------------------------------------------------------------------------
\usepackage{graphicx}
\usepackage{float}
\usepackage{caption}
\usepackage{subcaption}
\usepackage{xcolor}
\usepackage{tabularx}
\usepackage{longtable}
\usepackage{booktabs}
\usepackage{array}
\usepackage{enumitem}
\usepackage{hyperref}
\usepackage{fancyhdr}
\usepackage{titlesec}
\usepackage[most]{tcolorbox}
\usepackage{lastpage}

% ----------------------------------------------------------------------------
% Configuration des liens PDF
% ----------------------------------------------------------------------------
\hypersetup{
  colorlinks=true,
  linkcolor=black,
  urlcolor=blue,
  citecolor=black,
  pdftitle={[Titre du dossier professionnel]},
  pdfauthor={[Auteur]},
  pdfsubject={[Sujet]},
  pdfkeywords={[Mot-clé 1], [Mot-clé 2], [Mot-clé 3]}
}

% ----------------------------------------------------------------------------
% Couleurs du document
% ----------------------------------------------------------------------------
\definecolor{maincolor}{HTML}{1F2937}
\definecolor{accentcolor}{HTML}{2563EB}
\definecolor{softgray}{HTML}{F3F4F6}
\definecolor{midgray}{HTML}{6B7280}
\definecolor{lightborder}{HTML}{D1D5DB}
\definecolor{todoColor}{HTML}{B45309}
\definecolor{placeholderColor}{HTML}{374151}

% ----------------------------------------------------------------------------
% En-têtes et pieds de page
% ----------------------------------------------------------------------------
\pagestyle{fancy}
\fancyhf{}
\lhead{\small \documenttype}
\rhead{\small \projectname}
\cfoot{\small Page \thepage{} / \pageref{LastPage}}
\renewcommand{\headrulewidth}{0.4pt}
\renewcommand{\footrulewidth}{0.4pt}

% ----------------------------------------------------------------------------
% Style des titres
% ----------------------------------------------------------------------------
\titleformat{\chapter}
  {\normalfont\Huge\bfseries\color{maincolor}}
  {\thechapter.}{1em}{}
\titleformat{\section}
  {\normalfont\Large\bfseries\color{maincolor}}
  {\thesection}{1em}{}
\titleformat{\subsection}
  {\normalfont\large\bfseries\color{maincolor}}
  {\thesubsection}{1em}{}

% ----------------------------------------------------------------------------
% Variables principales du dossier
% ----------------------------------------------------------------------------
\newcommand{\documenttype}{Cahier des charges}
\newcommand{\projectname}{Psalmacor}
\newcommand{\authorname}{Daniel Medina}
\newcommand{\formationname}{Nom de la formation}
\newcommand{\organizationname}{Organisme / entreprise}
\newcommand{\versionnumber}{Version 1}
\newcommand{\documentdate}{Date}
\newcommand{\contactemail}{adresse@email.fr}

% ----------------------------------------------------------------------------
% Commandes réutilisables
% ----------------------------------------------------------------------------
\newcommand{\placeholder}[1]{\textcolor{placeholderColor}{\texttt{[#1]}}}
\newcommand{\todo}[1]{\textcolor{todoColor}{\textbf{À compléter :} #1}}
\newcommand{\keyword}[1]{\textbf{#1}}
\newcommand{\important}[1]{\textbf{\color{accentcolor}#1}}
\newcommand{\blankspace}[1]{\vspace*{#1}}

% Page vide structurée, prête à remplir
\newcommand{\pageapreenplir}[2]{%
  \section{#1}
  \begin{tcolorbox}[colback=softgray,colframe=lightborder,title=Zone à compléter]
    #2
  \end{tcolorbox}
  \blankspace{12cm}
}

% Bloc texte standard
\newtcolorbox{infobox}[1]{
  colback=softgray,
  colframe=accentcolor,
  title=#1,
  fonttitle=\bfseries
}

% Bloc décision
\newtcolorbox{decisionbox}[1]{
  colback=white,
  colframe=maincolor,
  title=Décision : #1,
  fonttitle=\bfseries
}

% Bloc exigence / besoin
\newtcolorbox{requirementbox}[1]{
  colback=white,
  colframe=accentcolor,
  title=Besoin / exigence : #1,
  fonttitle=\bfseries
}

% Image réelle avec légende
% Usage : \imageaveclegende[0.85\linewidth]{chemin/image.png}{Légende}{fig:mon-label}
\newcommand{\imageaveclegende}[4][0.85\linewidth]{%
  \begin{figure}[H]
    \centering
    \includegraphics[width=#1]{#2}
    \caption{#3}
    \label{#4}
  \end{figure}
}

% Emplacement image vide avec légende
% Usage : \imageplaceholder[0.85\linewidth]{6cm}{Légende}{fig:placeholder}
\newcommand{\imageplaceholder}[4][0.85\linewidth]{%
  \begin{figure}[H]
    \centering
    \fbox{%
      \begin{minipage}[c][#2][c]{#1}
        \centering
        \textcolor{midgray}{\Large Emplacement image}\\[0.4cm]
        \textcolor{midgray}{#3}
      \end{minipage}%
    }
    \caption{#3}
    \label{#4}
  \end{figure}
}

% Deux images côte à côte
% Usage : \deuximages{img1.png}{Légende 1}{fig:a}{img2.png}{Légende 2}{fig:b}
\newcommand{\deuximages}[6]{%
  \begin{figure}[H]
    \centering
    \begin{subfigure}{0.47\linewidth}
      \centering
      \includegraphics[width=\linewidth]{#1}
      \caption{#2}
      \label{#3}
    \end{subfigure}
    \hfill
    \begin{subfigure}{0.47\linewidth}
      \centering
      \includegraphics[width=\linewidth]{#4}
      \caption{#5}
      \label{#6}
    \end{subfigure}
  \end{figure}
}

% Tableau simple clé / valeur
\newenvironment{keyvaluetable}{%
  \begin{tabular}{>{\bfseries}p{0.30\linewidth}p{0.55\linewidth}}
}{%
  \end{tabular}
}

% Ligne de tableau clé / valeur
\newcommand{\kvrow}[2]{#1 & #2 \\\midrule}

% ----------------------------------------------------------------------------
% Début du document
% ----------------------------------------------------------------------------
\begin{document}

% ----------------------------------------------------------------------------
% Page de couverture
% ----------------------------------------------------------------------------
\begin{titlepage}
  \centering
  \vspace*{2cm}

  {\Huge\bfseries \documenttype \par}
  \vspace{1cm}
  {\LARGE\bfseries \projectname \par}

  \vfill

  \begin{tcolorbox}[colback=softgray,colframe=lightborder,width=0.85\linewidth]
    \begin{keyvaluetable}
      \kvrow{Auteur}{\authorname}
      \kvrow{Formation}{\formationname}
      \kvrow{Organisation}{\organizationname}
      \kvrow{Version}{\versionnumber}
      \kvrow{Date}{\documentdate}
      \kvrow{Contact}{\contactemail}
    \end{keyvaluetable}
  \end{tcolorbox}

  \vfill
  {\small Document préparé dans le cadre de \placeholder{contexte du dossier}.}
\end{titlepage}

% ----------------------------------------------------------------------------
% Historique des versions
% ----------------------------------------------------------------------------
\chapter*{Historique du document}
\addcontentsline{toc}{chapter}{Historique du document}

\begin{longtable}{p{0.15\linewidth}p{0.18\linewidth}p{0.22\linewidth}p{0.35\linewidth}}
\toprule
\textbf{Version} & \textbf{Date} & \textbf{Auteur} & \textbf{Modifications} \\
\midrule
\endhead
0.1 & \placeholder{Date} & \placeholder{Nom} & Création du document. \\
0.2 & \placeholder{Date} & \placeholder{Nom} & \placeholder{Description des modifications}. \\
\bottomrule
\end{longtable}

\newpage
\tableofcontents
\newpage
\listoffigures
\newpage
\listoftables
\newpage

% ============================================================================
% INTRODUCTION
% ============================================================================
\chapter{Introduction}

\section{Objet du document}
\begin{infobox}{But de cette section}
Présenter brièvement le rôle du dossier, son contexte de rédaction, son périmètre et les personnes concernées.
\end{infobox}

\todo{Décrire l'objet du dossier professionnel / cahier des charges.}
\blankspace{6cm}

\section{Présentation synthétique du projet}
\todo{Résumer le projet en quelques paragraphes : activité, service, site ou produit à réaliser, valeur proposée.}
\blankspace{8cm}

\section{Méthodologie de rédaction}
\todo{Indiquer la méthode suivie : entretiens, analyse de l'existant, étude concurrentielle, définition des besoins, arbitrages techniques.}
\blankspace{6cm}

% ============================================================================
% PRÉSENTATION DÉTAILLÉE
% ============================================================================
\chapter{Présentation détaillée}

\section{Présentation de l'activité}
\todo{Décrire l'activité, son origine, son positionnement, son vocabulaire, son univers et ses services principaux.}
\blankspace{8cm}

\section{Présentation du produit ou service}
\todo{Présenter le produit numérique ou le service à concevoir : site vitrine, espace de réservation, paiement, contenus, pages, tunnel utilisateur.}
\blankspace{8cm}

\section{Périmètre du projet}
\begin{requirementbox}{Périmètre fonctionnel}
Indiquer ce qui est inclus, ce qui est exclu et ce qui pourra être traité dans une version ultérieure.
\end{requirementbox}

\begin{itemize}[leftmargin=*]
  \item Inclus : \placeholder{élément inclus 1}
  \item Inclus : \placeholder{élément inclus 2}
  \item Exclu : \placeholder{élément exclu 1}
  \item Reporté : \placeholder{élément prévu pour une version future}
\end{itemize}

\blankspace{4cm}

\section{Intervenants}
\begin{longtable}{p{0.22\linewidth}p{0.22\linewidth}p{0.22\linewidth}p{0.24\linewidth}}
\toprule
\textbf{Nom} & \textbf{Rôle} & \textbf{Responsabilités} & \textbf{Contact} \\
\midrule
\endhead
\placeholder{Nom} & \placeholder{Client} & \placeholder{Validation, contenu, décisions} & \placeholder{Email / téléphone} \\
\placeholder{Nom} & \placeholder{Chef de projet} & \placeholder{Coordination, planning} & \placeholder{Email / téléphone} \\
\placeholder{Nom} & \placeholder{Développeur} & \placeholder{Conception, intégration, tests} & \placeholder{Email / téléphone} \\
\placeholder{Nom} & \placeholder{Référent pédagogique} & \placeholder{Suivi, validation pédagogique} & \placeholder{Email / téléphone} \\
\bottomrule
\caption{Liste des intervenants du projet.}
\label{tab:intervenants}
\end{longtable}

% ============================================================================
% CONTEXTE
% ============================================================================
\chapter{Contexte}

\section{Contexte général}
\todo{Décrire le contexte économique, social, culturel, pédagogique ou professionnel dans lequel le projet s'inscrit.}
\blankspace{8cm}

\section{Problématique}
\todo{Formuler le problème principal auquel le projet doit répondre.}
\blankspace{6cm}

\section{Enjeux}
\begin{itemize}[leftmargin=*]
  \item Enjeu commercial : \placeholder{description}
  \item Enjeu utilisateur : \placeholder{description}
  \item Enjeu technique : \placeholder{description}
  \item Enjeu organisationnel : \placeholder{description}
  \item Enjeu légal ou réglementaire : \placeholder{description}
\end{itemize}
\blankspace{4cm}

% ============================================================================
% OBJECTIFS
% ============================================================================
\chapter{Objectifs}

\section{Objectifs principaux}
\begin{longtable}{p{0.12\linewidth}p{0.46\linewidth}p{0.32\linewidth}}
\toprule
\textbf{ID} & \textbf{Objectif} & \textbf{Indicateur de réussite} \\
\midrule
\endhead
OBJ-01 & \placeholder{Objectif principal 1} & \placeholder{Indicateur mesurable} \\
OBJ-02 & \placeholder{Objectif principal 2} & \placeholder{Indicateur mesurable} \\
OBJ-03 & \placeholder{Objectif principal 3} & \placeholder{Indicateur mesurable} \\
\bottomrule
\caption{Objectifs principaux du projet.}
\label{tab:objectifs}
\end{longtable}

\section{Objectifs secondaires}
\todo{Lister les objectifs utiles mais non critiques.}
\blankspace{6cm}

\section{Critères de succès}
\todo{Définir les critères permettant de considérer le projet comme terminé et exploitable.}
\blankspace{6cm}

% ============================================================================
% CIBLE
% ============================================================================
\chapter{Cible}

\section{Public cible}
\todo{Décrire le ou les publics visés : âge, besoins, niveau numérique, attentes, freins, motivations.}
\blankspace{7cm}

\section{Personas}
\begin{longtable}{p{0.18\linewidth}p{0.22\linewidth}p{0.26\linewidth}p{0.24\linewidth}}
\toprule
\textbf{Persona} & \textbf{Profil} & \textbf{Besoins} & \textbf{Freins} \\
\midrule
\endhead
\placeholder{Persona 1} & \placeholder{Âge, contexte} & \placeholder{Besoins principaux} & \placeholder{Freins principaux} \\
\placeholder{Persona 2} & \placeholder{Âge, contexte} & \placeholder{Besoins principaux} & \placeholder{Freins principaux} \\
\placeholder{Persona 3} & \placeholder{Âge, contexte} & \placeholder{Besoins principaux} & \placeholder{Freins principaux} \\
\bottomrule
\caption{Personas représentatifs de la cible.}
\label{tab:personas}
\end{longtable}

\section{Parcours utilisateur type}
\todo{Décrire les étapes suivies par un utilisateur depuis la découverte du service jusqu'à l'action attendue.}
\blankspace{7cm}

% ============================================================================
% CONCURRENCE
% ============================================================================
\chapter{Concurrence}

\section{Identification des concurrents}
\begin{longtable}{p{0.22\linewidth}p{0.25\linewidth}p{0.25\linewidth}p{0.18\linewidth}}
\toprule
\textbf{Concurrent} & \textbf{Positionnement} & \textbf{Forces / faiblesses} & \textbf{Lien} \\
\midrule
\endhead
\placeholder{Concurrent 1} & \placeholder{Positionnement} & \placeholder{Forces et faiblesses} & \placeholder{URL} \\
\placeholder{Concurrent 2} & \placeholder{Positionnement} & \placeholder{Forces et faiblesses} & \placeholder{URL} \\
\placeholder{Concurrent 3} & \placeholder{Positionnement} & \placeholder{Forces et faiblesses} & \placeholder{URL} \\
\bottomrule
\caption{Analyse synthétique de la concurrence.}
\label{tab:concurrence}
\end{longtable}

\section{Benchmark visuel et fonctionnel}
\todo{Comparer les choix graphiques, les contenus, les fonctionnalités, les tunnels de conversion et les points de différenciation.}

\imageplaceholder{7cm}{Capture d'écran ou benchmark visuel à insérer}{fig:benchmark-placeholder}

\section{Positionnement différenciant}
\todo{Expliquer ce qui différencie le projet : ton, identité, offre, expérience, prix, accessibilité, proximité, spécialisation.}
\blankspace{6cm}

% ============================================================================
% ANALYSE MARKETING
% ============================================================================
\chapter{Analyse marketing}

\section{Proposition de valeur}
\begin{infobox}{Formulation de la proposition de valeur}
\placeholder{Pour [cible], le projet [nom] permet de [bénéfice principal] grâce à [différenciation].}
\end{infobox}
\blankspace{5cm}

\section{Analyse SWOT}
\begin{longtable}{p{0.23\linewidth}p{0.23\linewidth}p{0.23\linewidth}p{0.23\linewidth}}
\toprule
\textbf{Forces} & \textbf{Faiblesses} & \textbf{Opportunités} & \textbf{Menaces} \\
\midrule
\endhead
\placeholder{Force 1} & \placeholder{Faiblesse 1} & \placeholder{Opportunité 1} & \placeholder{Menace 1} \\
\placeholder{Force 2} & \placeholder{Faiblesse 2} & \placeholder{Opportunité 2} & \placeholder{Menace 2} \\
\placeholder{Force 3} & \placeholder{Faiblesse 3} & \placeholder{Opportunité 3} & \placeholder{Menace 3} \\
\bottomrule
\caption{Analyse SWOT du projet.}
\label{tab:swot}
\end{longtable}

\section{Stratégie de communication}
\todo{Décrire les canaux : référencement naturel, réseaux sociaux, newsletter, bouche-à-oreille, partenariats, contenus éditoriaux.}
\blankspace{6cm}

\section{Objectifs de conversion}
\begin{longtable}{p{0.18\linewidth}p{0.35\linewidth}p{0.35\linewidth}}
\toprule
\textbf{Action} & \textbf{Objectif} & \textbf{Mesure} \\
\midrule
\endhead
\placeholder{Contact} & \placeholder{Obtenir des demandes qualifiées} & \placeholder{Nombre de formulaires envoyés} \\
\placeholder{Réservation} & \placeholder{Convertir les visiteurs en clients} & \placeholder{Nombre de réservations} \\
\placeholder{Inscription} & \placeholder{Construire une audience} & \placeholder{Nombre d'inscriptions} \\
\bottomrule
\caption{Actions de conversion attendues.}
\label{tab:conversion}
\end{longtable}

% ============================================================================
% DÉFINITION DES BESOINS
% ============================================================================
\chapter{Définition des besoins}

\section{Étude de l'existant}
\todo{Décrire l'existant : supports actuels, outils utilisés, limites, contenus disponibles, contraintes déjà identifiées.}
\blankspace{8cm}

\section{Énoncé du besoin}
\begin{requirementbox}{Besoin général}
\placeholder{Le commanditaire souhaite disposer de [solution] afin de [objectif], pour [cible], dans le respect de [contraintes principales].}
\end{requirementbox}
\blankspace{5cm}

\section{Besoins fonctionnels}
\begin{longtable}{p{0.13\linewidth}p{0.38\linewidth}p{0.18\linewidth}p{0.21\linewidth}}
\toprule
\textbf{ID} & \textbf{Besoin} & \textbf{Priorité} & \textbf{Critère d'acceptation} \\
\midrule
\endhead
BF-01 & \placeholder{Afficher une page de présentation} & Must & \placeholder{Critère vérifiable} \\
BF-02 & \placeholder{Afficher une offre ou un service} & Must & \placeholder{Critère vérifiable} \\
BF-03 & \placeholder{Permettre la prise de contact} & Must & \placeholder{Critère vérifiable} \\
BF-04 & \placeholder{Permettre la réservation} & Should & \placeholder{Critère vérifiable} \\
BF-05 & \placeholder{Gérer des contenus administrables} & Could & \placeholder{Critère vérifiable} \\
\bottomrule
\caption{Liste des besoins fonctionnels.}
\label{tab:besoins-fonctionnels}
\end{longtable}

\section{Besoins non fonctionnels}
\begin{longtable}{p{0.13\linewidth}p{0.38\linewidth}p{0.18\linewidth}p{0.21\linewidth}}
\toprule
\textbf{ID} & \textbf{Besoin} & \textbf{Priorité} & \textbf{Critère d'acceptation} \\
\midrule
\endhead
BNF-01 & \placeholder{Performance} & Must & \placeholder{Temps de chargement cible} \\
BNF-02 & \placeholder{Accessibilité} & Must & \placeholder{Niveau attendu} \\
BNF-03 & \placeholder{Sécurité} & Must & \placeholder{Mesure attendue} \\
BNF-04 & \placeholder{Maintenabilité} & Should & \placeholder{Structure documentaire} \\
BNF-05 & \placeholder{Compatibilité mobile} & Must & \placeholder{Affichage responsive} \\
\bottomrule
\caption{Liste des besoins non fonctionnels.}
\label{tab:besoins-non-fonctionnels}
\end{longtable}

% ============================================================================
% FONCTIONS DU PRODUIT
% ============================================================================
\chapter{Les fonctions du produit}

\section{Fonctions principales}
\begin{longtable}{p{0.12\linewidth}p{0.30\linewidth}p{0.30\linewidth}p{0.18\linewidth}}
\toprule
\textbf{ID} & \textbf{Fonction} & \textbf{Description} & \textbf{Priorité} \\
\midrule
\endhead
F-01 & \placeholder{Page d'accueil} & \placeholder{Présenter l'activité et orienter l'utilisateur} & Must \\
F-02 & \placeholder{Pages services} & \placeholder{Présenter les offres disponibles} & Must \\
F-03 & \placeholder{Contact} & \placeholder{Permettre à l'utilisateur de contacter le prestataire} & Must \\
F-04 & \placeholder{Réservation} & \placeholder{Permettre la demande ou validation d'une réservation} & Should \\
F-05 & \placeholder{Administration} & \placeholder{Modifier certains contenus} & Could \\
\bottomrule
\caption{Fonctions principales du produit.}
\label{tab:fonctions-principales}
\end{longtable}

\section{User stories}
\begin{longtable}{p{0.12\linewidth}p{0.68\linewidth}p{0.12\linewidth}}
\toprule
\textbf{ID} & \textbf{User story} & \textbf{Priorité} \\
\midrule
\endhead
US-01 & En tant que \placeholder{type d'utilisateur}, je veux \placeholder{action}, afin de \placeholder{bénéfice}. & Must \\
US-02 & En tant que \placeholder{type d'utilisateur}, je veux \placeholder{action}, afin de \placeholder{bénéfice}. & Should \\
US-03 & En tant que \placeholder{type d'utilisateur}, je veux \placeholder{action}, afin de \placeholder{bénéfice}. & Could \\
\bottomrule
\caption{User stories du projet.}
\label{tab:user-stories}
\end{longtable}

\section{Arborescence prévue}
\todo{Décrire les pages et leur organisation.}
\begin{itemize}[leftmargin=*]
  \item Accueil
  \item Présentation
  \item Services
  \begin{itemize}
    \item \placeholder{Service 1}
    \item \placeholder{Service 2}
    \item \placeholder{Service 3}
  \end{itemize}
  \item À propos / intervenants
  \item Contact
  \item Mentions légales
  \item Politique de confidentialité
\end{itemize}
\blankspace{4cm}

\section{Maquettes et zoning}
\imageplaceholder{8cm}{Zoning de la page d'accueil}{fig:zoning-accueil}
\imageplaceholder{8cm}{Maquette ou wireframe d'une page service}{fig:wireframe-service}

% ============================================================================
% CONTRAINTES
% ============================================================================
\chapter{Les contraintes}

\section{Contraintes techniques}
\begin{longtable}{p{0.18\linewidth}p{0.38\linewidth}p{0.34\linewidth}}
\toprule
\textbf{Catégorie} & \textbf{Contrainte} & \textbf{Impact / solution prévue} \\
\midrule
\endhead
Hébergement & \placeholder{Hébergement, nom de domaine, SSL} & \placeholder{Impact et choix technique} \\
Compatibilité & \placeholder{Navigateurs, mobile, tablette} & \placeholder{Responsive design et tests} \\
Performance & \placeholder{Temps de chargement attendu} & \placeholder{Optimisation images, cache} \\
Sécurité & \placeholder{Formulaire, données, accès} & \placeholder{Validation, HTTPS, protections} \\
Maintenance & \placeholder{Mises à jour, documentation} & \placeholder{Procédure prévue} \\
\bottomrule
\caption{Contraintes techniques.}
\label{tab:contraintes-techniques}
\end{longtable}

\section{Contraintes légales et réglementaires}
\begin{longtable}{p{0.24\linewidth}p{0.36\linewidth}p{0.30\linewidth}}
\toprule
\textbf{Sujet} & \textbf{Obligation / point de vigilance} & \textbf{Action prévue} \\
\midrule
\endhead
Mentions légales & \placeholder{Informations éditeur, hébergeur} & \placeholder{Page dédiée} \\
RGPD & \placeholder{Données collectées, finalité, durée} & \placeholder{Politique de confidentialité} \\
Cookies & \placeholder{Cookies strictement nécessaires ou mesure d'audience} & \placeholder{Bandeau ou information} \\
Accessibilité & \placeholder{Lisibilité, navigation clavier, alternatives textuelles} & \placeholder{Bonnes pratiques WCAG} \\
Droit d'auteur & \placeholder{Images, textes, polices, contenus} & \placeholder{Licences et crédits} \\
\bottomrule
\caption{Contraintes légales et réglementaires.}
\label{tab:contraintes-legales}
\end{longtable}

\section{Contraintes de coûts}
\todo{Indiquer les limites budgétaires, les coûts prévus, les arbitrages et les coûts récurrents.}
\blankspace{6cm}

\section{Contraintes de délais}
\todo{Indiquer les échéances, jalons, dates de validation, dépendances et marges de sécurité.}
\blankspace{6cm}

\section{Contraintes graphiques et éditoriales}
\todo{Décrire l'identité visuelle, le ton éditorial, les couleurs, les typographies, les images et les contenus à fournir.}
\blankspace{6cm}

% ============================================================================
% DÉROULEMENT DU PROJET
% ============================================================================
\chapter{Déroulement du projet}

\section{Organisation générale}
\todo{Décrire les étapes du projet, les méthodes de validation, les outils de suivi et les règles d'échange.}
\blankspace{7cm}

\section{Phasage}
\begin{longtable}{p{0.12\linewidth}p{0.28\linewidth}p{0.38\linewidth}p{0.12\linewidth}}
\toprule
\textbf{Phase} & \textbf{Nom} & \textbf{Contenu} & \textbf{Statut} \\
\midrule
\endhead
1 & Analyse & \placeholder{Recueil des besoins, analyse concurrence, cible} & \placeholder{À faire} \\
2 & Conception & \placeholder{Arborescence, wireframes, choix techniques} & \placeholder{À faire} \\
3 & Réalisation & \placeholder{Développement, intégration, contenus} & \placeholder{À faire} \\
4 & Tests & \placeholder{Tests fonctionnels, responsive, accessibilité} & \placeholder{À faire} \\
5 & Livraison & \placeholder{Mise en ligne, documentation, recette} & \placeholder{À faire} \\
\bottomrule
\caption{Phases du projet.}
\label{tab:phasage}
\end{longtable}

\section{Planning prévisionnel}
\begin{longtable}{p{0.18\linewidth}p{0.18\linewidth}p{0.18\linewidth}p{0.34\linewidth}}
\toprule
\textbf{Tâche} & \textbf{Début} & \textbf{Fin} & \textbf{Livrable ou validation} \\
\midrule
\endhead
\placeholder{Analyse} & \placeholder{Date} & \placeholder{Date} & \placeholder{Document d'analyse} \\
\placeholder{Maquettes} & \placeholder{Date} & \placeholder{Date} & \placeholder{Wireframes validés} \\
\placeholder{Développement} & \placeholder{Date} & \placeholder{Date} & \placeholder{Version testable} \\
\placeholder{Tests} & \placeholder{Date} & \placeholder{Date} & \placeholder{Recette} \\
\placeholder{Livraison} & \placeholder{Date} & \placeholder{Date} & \placeholder{Mise en production} \\
\bottomrule
\caption{Planning prévisionnel.}
\label{tab:planning}
\end{longtable}

\section{Gestion des risques}
\begin{longtable}{p{0.22\linewidth}p{0.18\linewidth}p{0.18\linewidth}p{0.32\linewidth}}
\toprule
\textbf{Risque} & \textbf{Probabilité} & \textbf{Impact} & \textbf{Mesure prévue} \\
\midrule
\endhead
\placeholder{Retard contenu} & \placeholder{Moyenne} & \placeholder{Fort} & \placeholder{Prévoir une date limite de fourniture} \\
\placeholder{Changement de périmètre} & \placeholder{Moyenne} & \placeholder{Moyen} & \placeholder{Valider les demandes par écrit} \\
\placeholder{Problème technique} & \placeholder{Faible} & \placeholder{Fort} & \placeholder{Prévoir solution alternative} \\
\bottomrule
\caption{Matrice simplifiée des risques.}
\label{tab:risques}
\end{longtable}

% ============================================================================
% LIVRABLES
% ============================================================================
\chapter{Livrables}

\section{Liste des livrables}
\begin{longtable}{p{0.14\linewidth}p{0.34\linewidth}p{0.20\linewidth}p{0.22\linewidth}}
\toprule
\textbf{ID} & \textbf{Livrable} & \textbf{Format} & \textbf{Critère de validation} \\
\midrule
\endhead
LIV-01 & \placeholder{Cahier des charges} & PDF / TEX & \placeholder{Validé par le commanditaire} \\
LIV-02 & \placeholder{Maquettes} & PNG / PDF / Figma & \placeholder{Maquettes validées} \\
LIV-03 & \placeholder{Code source} & Dépôt Git & \placeholder{Code versionné} \\
LIV-04 & \placeholder{Site livré} & URL & \placeholder{Site accessible} \\
LIV-05 & \placeholder{Documentation} & PDF / Markdown & \placeholder{Procédure claire} \\
\bottomrule
\caption{Livrables attendus.}
\label{tab:livrables}
\end{longtable}

\section{Procédure de recette}
\todo{Décrire comment le projet sera testé et validé : liste des tests, retours, corrections, validation finale.}
\blankspace{7cm}

% ============================================================================
% DEVIS
% ============================================================================
\chapter{Devis}

\section{Hypothèses de chiffrage}
\todo{Indiquer les hypothèses retenues pour le devis : périmètre, nombre de pages, fonctionnalités, délais, exclusions.}
\blankspace{5cm}

\section{Estimation des coûts}
\begin{longtable}{p{0.42\linewidth}p{0.14\linewidth}p{0.16\linewidth}p{0.18\linewidth}}
\toprule
\textbf{Poste} & \textbf{Quantité} & \textbf{Prix unitaire} & \textbf{Total} \\
\midrule
\endhead
Analyse et cadrage & \placeholder{X h} & \placeholder{X €} & \placeholder{X €} \\
Conception UX / UI & \placeholder{X h} & \placeholder{X €} & \placeholder{X €} \\
Développement front-end & \placeholder{X h} & \placeholder{X €} & \placeholder{X €} \\
Développement back-end & \placeholder{X h} & \placeholder{X €} & \placeholder{X €} \\
Intégration contenus & \placeholder{X h} & \placeholder{X €} & \placeholder{X €} \\
Tests et corrections & \placeholder{X h} & \placeholder{X €} & \placeholder{X €} \\
Mise en ligne & \placeholder{X h} & \placeholder{X €} & \placeholder{X €} \\
Maintenance optionnelle & \placeholder{X mois} & \placeholder{X €} & \placeholder{X €} \\
\midrule
\textbf{Total HT} & & & \placeholder{X €} \\
\textbf{TVA} & & & \placeholder{X €} \\
\textbf{Total TTC} & & & \placeholder{X €} \\
\bottomrule
\caption{Estimation budgétaire du projet.}
\label{tab:devis}
\end{longtable}

\section{Coûts récurrents}
\begin{longtable}{p{0.36\linewidth}p{0.22\linewidth}p{0.32\linewidth}}
\toprule
\textbf{Poste} & \textbf{Coût estimé} & \textbf{Fréquence} \\
\midrule
\endhead
Nom de domaine & \placeholder{X €} & Annuel \\
Hébergement & \placeholder{X €} & Mensuel / annuel \\
Outils tiers & \placeholder{X €} & Mensuel / annuel \\
Maintenance & \placeholder{X €} & Mensuel / ponctuel \\
\bottomrule
\caption{Coûts récurrents prévisionnels.}
\label{tab:couts-recurrents}
\end{longtable}

% ============================================================================
% ANNEXES
% ============================================================================
\appendix
\chapter{Annexes}

\section{Glossaire}
\begin{longtable}{p{0.26\linewidth}p{0.64\linewidth}}
\toprule
\textbf{Terme} & \textbf{Définition} \\
\midrule
\endhead
\placeholder{Terme 1} & \placeholder{Définition} \\
\placeholder{Terme 2} & \placeholder{Définition} \\
\placeholder{Terme 3} & \placeholder{Définition} \\
\bottomrule
\caption{Glossaire du projet.}
\label{tab:glossaire}
\end{longtable}

\section{Sources et références}
\begin{longtable}{p{0.32\linewidth}p{0.46\linewidth}p{0.12\linewidth}}
\toprule
\textbf{Source} & \textbf{Utilisation} & \textbf{Date} \\
\midrule
\endhead
\placeholder{Nom de la source} & \placeholder{Pourquoi cette source est utilisée} & \placeholder{Date} \\
\placeholder{Nom de la source} & \placeholder{Pourquoi cette source est utilisée} & \placeholder{Date} \\
\bottomrule
\caption{Sources et références utilisées.}
\label{tab:sources}
\end{longtable}

\section{Captures d'écran}
\imageplaceholder{7cm}{Capture d'écran 1 à insérer}{fig:annexe-capture-1}
\imageplaceholder{7cm}{Capture d'écran 2 à insérer}{fig:annexe-capture-2}

\section{Exemples de blocs réutilisables}

\begin{infobox}{Exemple de bloc d'information}
Ce bloc peut servir à présenter une précision, une hypothèse, une synthèse ou un rappel important.
\end{infobox}

\begin{decisionbox}{Nom de la décision}
\textbf{Contexte :} \placeholder{Contexte de la décision.}\\
\textbf{Décision retenue :} \placeholder{Choix retenu.}\\
\textbf{Justification :} \placeholder{Raison du choix.}
\end{decisionbox}

\begin{requirementbox}{Identifiant du besoin}
\textbf{Description :} \placeholder{Décrire le besoin.}\\
\textbf{Priorité :} \placeholder{Must / Should / Could / Won't.}\\
\textbf{Critère d'acceptation :} \placeholder{Condition vérifiable.}
\end{requirementbox}

\section{Page libre}
\todo{Utiliser cette page pour ajouter des schémas, notes, captures, comptes rendus ou éléments complémentaires.}
\blankspace{14cm}

\end{document}
